\section{Notations}\label{sec:notations}

To maintain uniformity, this document employs a set of common mathematical notations. Vectors are represented as bold lower-case characters, such as~\(\bm{v}\). Matrices are represented as bold upper-case characters, such as~\(\bm{M}\). Digits printed in bold represent vectors of constant values, such as~\(\bm{1}={[1, 1, \ldots, 1]}^\top\).

Let~\(\bm{M} = [m_{ij}] \in \R^{m\times n}\) denote a matrix with entry~\(m_{ij}\) at row~\(i\), column~\(j\), having~\(i\)-th row~\(\bm{m}_{i,:}\) and~\(j\)-th column~\(\bm{m}_{:,j}\). A vector-of-vectors (VoV) of~\({\{v_{i}\}}_{1:l}\) is a vector~\(\bm{v}\in\Omega^n\) composed of sub-vectors~\(\bm{v}_i\in\Omega^{n_i}\), where~\(\bm{v}={[\bm{v}_1^\top, \bm{v}_2^\top, \ldots, \bm{v}_l^\top]}^\top\) and~\(\sum_{i=1}^l n_i = n\).

A scalar component of a VoV is~\({[\bm{v}_i]}_j\). For a given VoV,~\(\bm{v}\), the following set notations apply.
\begin{equation*}
  \bm{v}_{i\neq j} =
  \begin{cases}
    {\left[
        \bm{v}_2^\top,
        \bm{v}_3^\top,
        \cdots,
        \bm{v}_l^\top
    \right]}^\top & \text{for }\ j = 1\\
    %
    {\left[
        \bm{v}_1^\top,
        \bm{v}_2^\top,
        \cdots,
        \bm{v}_{l-1}^\top
    \right]}^\top & \text{for }\ j = l\\
    %
    {\left[
        \bm{v}_1^\top,
        \cdots,
        \bm{v}_{j-1}^\top,
        \bm{v}_{j+1}^\top,
        \cdots,
        \bm{v}_l^\top
    \right]}^\top & \text{for any }\ j\in\{2,\ldots,l-1\}
  \end{cases}
\end{equation*}
\begin{alignat*}{3}
  &\bm{v}_{i:j} &&= {\left[
      \bm{v}_i^\top,
      \bm{v}_{i+1}^\top,
      \cdots,
      \bm{v}_{j-1}^\top,
      \bm{v}_j^\top
  \right]}^\top\quad &&\text{for }\ 1\leq i<j \leq l\\
  %
  &{\left[\bm{v}_k\right]}_{i:j} &&= {\left[
      {\left[\bm{v}_k\right]}_i,
      {\left[\bm{v}_k\right]}_{i+1},
      \cdots,
      {\left[\bm{v}_k\right]}_{j-1},
      {\left[\bm{v}_k\right]}_j
  \right]}^\top\quad &&\text{for }\ 1\leq i<j \leq n_k
\end{alignat*}

A scalar component of a regular vector retains the conventional notation~\(v_i\). The~\(n^\text{th}\) norm of a vector is formulated as~\(\|\bm{v}\|_n = {(\sum_{i=1}^n |v_i|^n)}^{1/n}\). For the second-order (Euclidean) norm, the subscript 2 is omitted, written as~\(\|\bm{v}\|\).

For optimization processes, the following conventions are used. Variables calculated at the system level are denoted with an overline (e.g.,~\(\overline{\bm{v}}\)), while those calculated at the subsystem level are denoted with an underline (e.g.,~\(\underline{\bm{v}}\)). The state of a variable,~\(\bm{v}\), at iteration instance~\(k\) is denoted as~\(\bm{v}^{(k)}\). An asterisk superscript,~\(\bm{v}^*\), marks the best solution found found.

Finally, function mapping is generalized. If~\(f: \R \to \R\), then~\(f(\bm{x} \in \R^n)\) implies element-wise application resulting in a vector in~\(\R^n\). Similarly, if~\(f: \R^n \to \R\), then~\(f(\bm{X} \in \R^{n \times m})\) implies application to each column of matrix~\(\bm{X}\), resulting in a vector in~\(\R^m\).
